\begin{textsample}{POS Dim 6 – rn – Score 6.00 – rn/rn\_em\_44.txt}  \label{ex:f6_pos_237}
7.
DIRETRIZES PARA OS ITINERÁRIOS \textbf{FORMATIVOS} A Arquitetura Curricular do Ensino Médio Potiguar é composta por Formação Geral Básica, que compreende as aprendizagens essenciais, como apresentado no capítulo anterior deste Referencial, e os Itinerários \textbf{Formativos}, enquanto estratégia de flexibilização curricular, que promovem o aprofundamento e ampliação das aprendizagens, visando a articulação da instituição escolar aos anseios da comunidade em que está inserida e o fomento ao protagonismo estudantil.
Os Itinerários \textbf{Formativos}, portanto, são uma oportunidade de diversificação da oferta, promovendo interação entre áreas do conhecimento e atendendo às demandas dos jovens e da comunidade.
A Lei nº 12.796/2013, que alterou a Lei de Diretrizes e Bases da Educação Nacional - LDB ( Lei nº 9.394/1996 ), em seu Art. 26 preconiza : > Os currículos da educação infantil, do ensino fundamental e do ensino médio devem ter base nacional comum, a ser complementada, em cada sistema de ensino e em cada estabelecimento escolar, por uma parte diversificada, exigida pelas características regionais e locais da sociedade, da cultura, da economia e dos educandos. ( Art. 26, LDB ).
Corroborando com essa perspectiva, o Plano Nacional de Educação - PNE ( Lei nº 13.005/2014 ) prevê na Estratégia 3.1 da Meta 3 a construção de currículos escolares que se organizem de forma flexível e diversificada e por meio de partes obrigatória e eletiva.
Ainda, o Plano Estadual de Educação - PEE ( Lei nº 10.049/2016 ) também aponta para uma flexibilização curricular quando em sua Estratégia 3 da Meta 3 é ratificada a permanência dos jovens na escola e, consequentemente, a conclusão da Etapa do Ensino Médio ao conectar o currículo escolar às necessidades e expectativas \textbf{formativas} dos estudantes das diversas comunidades escolares no território potiguar.
Os Itinerários \textbf{Formativos}, nesse sentido, estão articulados indissociavelmente à Formação Geral Básica, buscando a ampliação das aprendizagens ora desenvolvidas.
De acordo com as Diretrizes Curriculares Nacionais do Ensino Médio ( Resolução CEB nº 3/2018 ), os Itinerários \textbf{formativos} são definidos como : > cada conjunto de unidades curriculares ofertadas pelas instituições e redes de ensino que possibilitam ao estudante aprofundar seus conhecimentos e se preparar para o prosseguimento de estudos ou para o mundo do trabalho de forma a contribuir para a construção de soluções de \textbf{problemas} específicos da sociedade ( Art. 6º, inciso III, DCNEM/2018 ).
Norteados por Eixos \textbf{Estruturantes}, os Itinerários \textbf{Formativos} exercem papel fundamental na flexibilização curricular, ofertando práticas pedagógicas interdisciplinares e transdisciplinares.
De acordo com o disposto nos incisos do § 2º, do Art. 12 das DCNEM/2018, os Eixos são complementares entre si e orientam os Itinerários \textbf{Formativos} na ampliação e aprofundamento das aprendizagens.
São eles : investigação científica, processos criativos, \textbf{mediação} e intervenção sociocultural e \textbf{empreendedorismo}.
As condições para flexibilização são garantidas na construção do currículo escolar, em um processo dialógico com a comunidade, oportunizando ao estudante a participação desde o processo de concepção até o momento de escolha do Itinerário \textbf{Formativo}.
A construção do trabalho pedagógico, com engajamento da comunidade escolar, propicia ao estudante o protagonismo no seu percurso \textbf{formativo}, na medida em que este se percebe como indivíduo e ao outro na sua diversidade e coletividade, entende o meio em que está inserido, as condições de vida e as múltiplas possibilidades de intervenção e transformação da realidade.
No contexto da flexibilização, este Referencial Curricular apresenta o Projeto de Vida como uma unidade curricular fixa que promove a reflexão nas dimensões pessoal, cidadã e profissional, apoiando o estudante na escolha do percurso \textbf{formativo} e, sobretudo, das \textbf{Trilhas} de Aprofundamento ofertadas.
O Projeto de Vida se constitui, portanto, como prática transversal ao trabalho pedagógico desenvolvido no ambiente escolar.

% matched lemmas: empreendedorismo, estruturante, formativo, mediação, problema, trilha
\end{textsample}
